%% ma: L11-B16-0087
%% lop: 11
%% chuong: 5
%% bai: 16
%% dang: 11.16.1
%% loai: TL
%% muc: VD
%% dap_an: "-"
%% nguon: "Ảnh giáo viên gửi 25/11/2026: Giới hạn hàm số dạng 0/0 (3 bài, có ĐS), Bài 2 câu 5"
%% kien_thuc: "Lớp 11 Bài 16 (giới hạn của hàm số tại một điểm); chia đa thức theo sơ đồ Horner, phân tích đa thức thành nhân tử"
%% trang_thai: chinh-thuc
\begin{ex}
Tính giới hạn $A=\lim\limits_{x\to-\sqrt3}\dfrac{2x^3-3x^2+x+9+7\sqrt3}{3-x^2}$.
\loigiai{Chia cho $x+\sqrt3$: $2x^3-3x^2+x+9+7\sqrt3=(x+\sqrt3)\left[2x^2-(3+2\sqrt3)x+7+3\sqrt3\right]$; $3-x^2=(\sqrt3-x)(\sqrt3+x)$. Vậy $A=\dfrac{6+(3+2\sqrt3)\sqrt3+7+3\sqrt3}{2\sqrt3}=\dfrac{19+6\sqrt3}{2\sqrt3}=\dfrac{18+19\sqrt3}{6}$.}
\end{ex}
